\section{Related Work}

We situate \axvisor{} against three bodies of work: hypervisor architectures, OS
modularity and library-OS designs, and software-architecture approaches to
portability.

\paragraph{Hypervisor architectures and their substrate coupling}
Xen introduced paravirtualization to let multiple commodity OSes share hardware
with low overhead, but its design presupposes a privileged Dom0 and a specific
hypervisor--guest contract~\cite{barham2003xen}. KVM turned the Linux kernel
itself into a hypervisor via a character device and ioctl
ABI~\cite{kivity2007kvm}, which made virtualization a Linux subsystem rather
than a portable component; Firecracker builds on KVM and inherits that
dependency while specializing for serverless density~\cite{agache2020firecracker}.
Static-partitioning hypervisors reduce coupling to a substrate but do not
eliminate it: Jailhouse relies on a booted Linux to initialize hardware before
partitioning~\cite{ramsauer2017jailhouse}; Bao assumes a bare-metal boot and
memory model~\cite{martins2020bao}; ACRN couples to a service-VM configuration
for IoT consolidation~\cite{li2019acrn}. Comparative studies of these systems
treat each as an independent whole and do not ask whether a shared
virtualization core could span substrates~\cite{lozano2023comprehensive}.
Research designs that rethink the boundary---DuVisor's user-level hypervisor
through delegated virtualization~\cite{chen2022duvisor} and the
Out-of-Hypervisor approach to nested virtualization~\cite{bitchebe2022ooh}---
relocate coupling between kernel and user space rather than abstracting the
substrate contract. \axvisor{} differs by making the substrate dependency an
explicit, minimal contract that multiple OSes implement, so the same core is
reused rather than re-embedded.

\paragraph{OS modularity, library OSes, and unikernels}
Unikraft decomposes OS construction into selectable primitives so that a
specialized unikernel can be assembled per application, demonstrating that OS
services can be factored into reusable libraries~\cite{kuenzer2021unikraft}.
LibrettOS fuses library-OS and multiserver paradigms to balance isolation,
performance, and adaptability~\cite{nikolaev2020librettos}. Unikernel isolation
continues to be explored for edge and embedded targets where a small footprint
and fast boot matter~\cite{kaiser2024unikernels, ueda2024mewz}. Asterinas shows
that a Rust-based framekernel can preserve Linux ABI compatibility while keeping
a small, sound trusted computing base~\cite{peng2025asterinas,
peng2024framekernel}. These works establish that OS services can be modular and
substrate-agnostic in principle; \axvisor{} applies the same principle from the
consumer side, defining what a hypervisor needs from such a substrate and
showing that the requirement can be met by unikernel, framekernel, and
monolithic OSes alike.

\paragraph{Guest-facing standardization and hardware portability}
The KVM ABI itself is a form of guest-side decoupling: by standardizing the
VMM-to-hypervisor interface, it lets guests and toolchains be reused across KVM
implementations~\cite{kivity2007kvm}. \axvisor{} adopts this ABI precisely to
inherit that ecosystem on the guest side while adding substrate-side decoupling
below. Separately, hardware portability across ISAs is an active axis: work on
Arm and RISC-V virtualization explores how hypervisors span
architectures~\cite{sa2023cva6}. Hardware portability is orthogonal to
substrate-OS portability---a hypervisor can be cross-ISA yet bound to one OS, or
cross-OS yet single-ISA---and \axvisor{}'s contribution complements rather than
competes with this axis. Work on virtualization for safety-critical systems,
including virtio-based real-time I/O sharing~\cite{peixoto2024birtio} and
timing-impossibility results for perfect hypervisors~\cite{guri2025impossibility},
addresses different concerns (predictability and fundamental limits) and is
complementary to portability.

\paragraph{Software architecture and portability}
Abstraction and modularity are foundational to managing complexity and enabling
reuse in software systems~\cite{shaw2025abstractions, wan2023software}. Yet
portability is often an aspiration that breaks down when abstractions leak
platform semantics~\cite{mince2023illusion}. \axvisor{} engages this tension
directly: its substrate contract is designed to surface, not conceal, the
semantic differences (physical memory ownership, interrupt context, stage-2
control) where portability is contested, and its evaluation measures adaptation
cost as evidence that the abstraction holds.
